\documentclass[conference]{IEEEtran}
\IEEEoverridecommandlockouts

\usepackage{amsmath,amssymb,amsfonts}
\usepackage{amsthm}
\usepackage{graphicx}
\usepackage{textcomp}
\usepackage{xcolor}
\usepackage{booktabs}
\usepackage{tabularx}
\usepackage{makecell}
\usepackage{adjustbox}

\usepackage[style=numeric, backend=biber]{biblatex}
\addbibresource{references.bib}

\usepackage{algorithm}
\usepackage{algpseudocode}

\theoremstyle{definition}
\newtheorem{definition}{Definition}

\theoremstyle{plain}
\newtheorem{theorem}{Theorem}
\newtheorem{lemma}{Lemma}
\newtheorem{proposition}{Proposition}
\newtheorem{corollary}{Corollary}

\usepackage{microtype}
\usepackage{listings}

\def\BibTeX{{\rm B\kern-.05em{\sc i\kern-.025em b}\kern-.08em
    T\kern-.1667em\lower.7ex\hbox{E}\kern-.125emX}}

\definecolor{codegray}{rgb}{0.5,0.5,0.5}
\definecolor{codeback}{rgb}{0.97,0.97,0.97}
\lstset{
    backgroundcolor=\color{codeback},
    basicstyle=\ttfamily\small,
    breaklines=true,
    captionpos=b,
    commentstyle=\color{codegray},
    frame=single,
    keepspaces=true,
    numbers=none,
    keywordstyle=\color{black}\bfseries,
    showspaces=false,
    showstringspaces=false,
    showtabs=false,
    tabsize=4
}

\usepackage{hyperref}

\begin{document}

\title{Beyond MDS: A Generation-Order Defense Closing Three Independent\\
2026 Attacks on Poseidon's Linear Layer}

\author{
\IEEEauthorblockN{Chaolu Liu}
\IEEEauthorblockA{\textit{School of Computer Science} \\
\textit{University College Dublin}\\
Dublin, Ireland \\
cl5248@hotmail.com}
\and
\IEEEauthorblockN{Emilio Martinez}
\IEEEauthorblockA{\textit{School of Computing} \\
\textit{National College of Ireland}\\
Dublin, Ireland \\
eyolticmg@icloud.com}
\and
\IEEEauthorblockN{Furqan Rustam}
\IEEEauthorblockA{\textit{School of Computer Science} \\
\textit{University College Dublin}\\
Dublin, Ireland \\
furqan.rustam1@gmail.com}
\and
\IEEEauthorblockN{Anca Delia Jurcut}
\IEEEauthorblockA{\textit{School of Computer Science} \\
\textit{University College Dublin}\\
Dublin, Ireland \\
anca.jurcut@ucd.ie}
}

\maketitle

\begin{abstract}
Poseidon is among the most widely deployed arithmetization-oriented
permutations in zero-knowledge proof systems, and its security rests on
a single structural assumption: that the linear layer mixing its partial
rounds is maximum distance separable (MDS). We show this assumption is
necessary but not sufficient. Three independent 2026 results --
Vitto's Slipway~\cite{slipway2026}, Merz and Rodr\'iguez Garc\'ia's
Skipping Class~\cite{cryptoeprint:2026/306}, and Jo's Midpoint
Reset~\cite{jo2026} -- each construct a properly MDS linear layer that
passes every published subspace-trail check and still breaks the
permutation's differential or collision security, through three
structurally distinct mechanisms. We identify the single assumption
common to all three: that the linear layer and the round constants are
generated independently of one another. We give (i) a unified structural
criterion generalizing both the GRS21 invariant-subspace
check~\cite{weaklayers2020} and the tensor-eigenvector check implicit in
Merz and Rodr\'iguez Garc\'ia's attack; (ii) a single, essentially-free
generation-order countermeasure that enforces the required independence
by construction; and (iii) the first independent replication of Jo's
full-round compression collision and the first empirical test of whether
this countermeasure defeats it -- across nine independent trials, zero
exceptions. We instantiate the result at a lightweight KoalaBear
parameter set, cross-validate it against a real \texttt{circom}/Groth16
compile and the production \texttt{circomlib} reference implementation,
subject it to an eight-point exhaustive falsification battery, audit
three real-world Poseidon deployments (\texttt{circomlib}, Plonky3,
\texttt{arkworks-rs}), and position its cost against Poseidon2, Anemoi,
and Polocolo using a contemporaneous, independently authored
survey~\cite{sok2026}. We report five further design levers we tested
and found did not help, as negative results. Two attack classes -- S-box
skipping~\cite{sboxskip2026} and bivariate resultant attacks against
Poseidon's own linear layer~\cite{resultant2026} -- remain outside this
paper's countermeasure and are stated explicitly as open problems.
\end{abstract}

\begin{IEEEkeywords}
Poseidon, Poseidon2, arithmetization-oriented primitives, algebraic
cryptanalysis, subspace trails, generation-order attacks, zero-knowledge
proofs, lightweight circuit design
\end{IEEEkeywords}

\section{Introduction}
\label{sec:introduction}

Arithmetization-oriented (AO) hash functions are a load-bearing
component of modern zero-knowledge proof systems: because their native
cost is measured in R1CS, AIR, or PLONK constraints rather than CPU
cycles, their internal structure is far more exposed to algebraic
attack than a classical hash function's. Poseidon~\cite{poseidon2021}
and its successor Poseidon2~\cite{poseidon22023}, built on the HADES
strategy~\cite{hades2020}, are among the most widely deployed primitives
in this category, used as the hash function inside Merkle trees,
commitment schemes, and Fiat--Shamir transcripts across production
zero-knowledge systems.

Poseidon's efficiency comes from \emph{partial rounds}: a nonlinear
S-box is applied to a single state coordinate, and a fixed linear layer
$M$ is relied upon to diffuse that nonlinearity across the remaining
coordinates before the next round. The folklore security requirement on
$M$ is that it be maximum distance separable (MDS). We confirm, as a
self-contained baseline (\S\ref{sec:background}), that this requirement
is \emph{necessary}: a non-MDS $M$ admits an explicit invariant subspace
in which the S-box never activates, collapsing the entire partial-round
sequence to a linear map. What this paper is about is that MDS is
\textbf{not sufficient}.

\paragraph{The problem.} Over the course of 2026, three independent
research groups -- developed in parallel, none derivative of the
others -- each constructed a linear layer that is properly MDS, passes
every subspace-trail check published to date, and still breaks a
standard Poseidon instance:

\begin{itemize}[leftmargin=1.2em]
  \item \textbf{Vitto}~\cite{slipway2026} shows every linear layer
  unavoidably admits \emph{finite} subspace trails, and that an
  adversary who chooses $M$ \emph{after} the round constants are fixed
  can steer a constrained input family through the initial full rounds
  into one, collapsing the effective algebraic degree far below the
  naive bound.
  \item \textbf{Merz and Rodr\'iguez Garc\'ia}~\cite{cryptoeprint:2026/306}
  show Poseidon2(b)'s efficient Kronecker-structured external matrix
  $M_\varepsilon=P_{t/4}\otimes M_4$ has an unconstrained
  $(q,-q,0,\ldots,0)$ eigenvector subspace that survives full rounds at
  zero algebraic-degree cost, independently of the vetted internal
  matrix.
  \item \textbf{Jo}~\cite{jo2026} shows that if the round constants are
  fixed before $M$ is chosen, an adversary can program an \emph{exact}
  nonlinear two-execution trajectory through every round via a linear
  interpolation $M=VU^{-1}$, yielding a genuine full-round
  \textbf{compression collision} -- not a distinguisher -- in a matrix
  that is MDS and clears every published subspace-trail check,
  including the Poseidon Cryptanalysis Initiative's own reference
  verifier.
\end{itemize}

Each paper proposes, at most, an ad hoc fix for the single mechanism it
exhibits. None states or validates a criterion covering more than one
mechanism. Jo's own discussion \emph{names}, without building or
testing, the fix we independently arrived at and validate in this paper.
A full read of the most comprehensive survey of this literature
available~\cite{sok2026} (134 pages, \S\ref{sec:sok}) confirms that no
existing source analyzes Jo's attack beyond a single combined sentence
shared with Vitto's, and performs no synthesis across any of the three
mechanisms above.

\paragraph{Our observation.} All three attacks succeed by violating one
assumption: \emph{the linear layer and the round constants must be
generated independently of one another}. Vitto's and Jo's constructions
both require the matrix to be chosen with knowledge of the constants (or
vice versa); Merz and Rodr\'iguez Garc\'ia's requires a specific,
unconstrained factor order that a standards process could equally well
avoid at zero extra cost. No static property of $M$ -- however
sophisticated the check -- can detect a violation of this assumption,
because the violation is a relation between $M$ and the constants, not a
property of $M$ in isolation. We confirm this is not merely a
theoretical gap: in \S\ref{sec:midpoint-reset}, our own replication of
Jo's construction yields a matrix that our strongest structural
certificate (\S\ref{sec:criterion}) reports fully secure.

\subsection{Contributions}
We state our contributions in three tiers of novelty rather than as one
undifferentiated list, since several components of this paper implement
or validate ideas proposed, but not built or tested, by others. Claiming
otherwise would misattribute prior work; we would rather a reviewer see
the distinction stated plainly than discover it themselves.

\paragraph{Tier 1 -- results we believe are genuinely new.}
\begin{enumerate}[leftmargin=1.3em]
  \item \textbf{The first independent replication of Jo's full-round
  collision} (\S\ref{sec:midpoint-reset}) \textbf{and the first empirical
  test of whether the countermeasure Jo's own discussion names, but
  neither builds nor tests, actually defeats it}: across nine independent
  trials (two matrix families, two toy fields, five seeds), zero
  instances of our binding being defeated. A full read of the most
  comprehensive available survey~\cite{sok2026} confirms no existing
  source analyzes Jo's attack beyond a one-sentence mention shared with
  Vitto's -- this is, to our knowledge, the only component of this paper
  that is unambiguously without prior treatment anywhere.
  \item \textbf{A Kronecker-agnostic generalization of the
  tensor-eigenvector check} (\S\ref{sec:criterion}, step (c) of
  Algorithm~\ref{alg:criterion}): Merz and Rodr\'iguez Garc\'ia's own
  finding~\cite{cryptoeprint:2026/306} is specific to a known Kronecker
  structure; our certificate independently rediscovers the identical
  vulnerable/safe distinction on the real $16\times16$ matrices
  \emph{without being given any tensor structure as input}, and we use
  it, unprompted, to find the same vulnerable factor order present in
  Plonky3's production code (\S\ref{sec:audit}).
  \item \textbf{Two concrete, previously undocumented real-world
  findings} (\S\ref{sec:audit}): \texttt{arkworks-rs}'s parameter
  generator asserts three well-formedness checks in source comments with
  zero corresponding code, and Plonky3's external matrix uses the
  mechanism's vulnerable factor order in an actively used STARK proving
  stack. Neither is a demonstrated break of a deployed instance; both
  are, as far as we are aware, newly documented.
  \item \textbf{Tensor-structuring the full-round matrix as an
  efficiency lever, not only a security fix} (\S\ref{sec:lightweight}):
  applying the safe Kronecker order to $M_{\mathrm{full}}$ (not just
  $M_I$) yields a measured $73\%$ reduction in real multiplicative
  complexity at $t=16$ -- a combination of the published security fix
  with an efficiency argument we did not find reported elsewhere.
\end{enumerate}

\paragraph{Tier 2 -- correct, validated engineering, built on others'
ideas.} \textbf{The generation-order countermeasure of
\S\ref{sec:countermeasure} is not our idea}: Vitto's~\cite{slipway2026}
\S6.3 and Jo's~\cite{jo2026} discussion both name binding round
constants to a canonical encoding of the matrix as the natural fix for
their respective attacks, and neither builds or tests it. Our
contribution is constructing it correctly (including finding and closing
a real gap, where an earlier version bound only $M_I$ and left
Vitto's actual target, $M_{\mathrm{full}}$, unbound) and being the first
to test it against Jo's strictly stronger attack (Tier 1, item 1).
Likewise, the GRS21 invariant-subspace criterion~\cite{weaklayers2020}
(Algorithm~\ref{alg:criterion}, steps (a)/(a-fast)/(b)) is 2020/2021
prior work; our contribution is a field-generic, independently validated
implementation, not a new criterion. The lightweight KoalaBear
instantiation, its cross-validation against \texttt{circomlib} and a
real Groth16 pipeline, and the eight-point exhaustive battery
(\S\ref{sec:lightweight}--\S\ref{sec:battery}) are rigorous validation
work establishing that our combination of known-good components behaves
as claimed; they are due diligence in support of the Tier 1 results, not
independent discoveries.

\paragraph{Tier 3 -- negative results, reported for completeness.}
\textbf{Five design levers we tested and found did not help}
(\S\ref{sec:levers}): lowering the S-box exponent, rotating the active
S-box coordinate, two tricks imported from GMiMCHash2, and a naive
lookup-hybrid degree model. None of these are contributions in the usual
sense -- we report them so the next attempt does not repeat the same
dead ends.

Finally, \textbf{an explicit scope boundary} (\S\ref{sec:discussion}):
two further 2026 attack classes remain open, confirmed as open by a full
read of the field's most comprehensive survey, not merely by our own
prior gap.

\section{Background}
\label{sec:background}

Let $\mathbb F_p$ be a prime field, $t$ the state width, $\alpha$ the
S-box exponent with $\gcd(\alpha,p-1)=1$. A \emph{full round} applies
$S_F(z)=z^{\circ\alpha}$ (componentwise power) to every coordinate; a
\emph{partial round} applies $S_P(z)=(z_0^{\alpha},z_1,\ldots,z_{t-1})^T$,
leaving all but the first coordinate linear. Each round is
\[
  R_r(x) = M\cdot\big(S_r(x) + c^{(r)}\big),
\]
S-box, then round-constant addition, then the linear layer $M$ -- the
convention used by this paper's own reference implementation. The full
permutation applies $R_F/2$ full rounds, $R_P$ partial rounds, then
$R_F/2$ more full rounds. Poseidon2(b)-style designs use two distinct,
stationary matrices: $M_I$ for the partial segment, $M_{\mathrm{full}}$
for the full-round prefix and suffix.

\begin{definition}[Branch number / MDS]
For $M\in\mathbb F_p^{t\times t}$,
$B_d(M)=\min_{\Delta\ne0}\big(\mathrm{wt}(\Delta)+\mathrm{wt}(M\Delta)\big)$.
$M$ is MDS iff $B_d(M)\ge t+1$.
\end{definition}

\begin{lemma}[Baseline: an explicit invariant subspace trail]
\label{lem:invariant-subspace}
Let $t=3$, $M$ non-MDS, with $\Delta=(0,0,\delta)\ne0$ and
$M\Delta \in V:=\{(0,0,x):x\in\mathbb F_p\}$. Then $V$ is invariant under
the full sequence of $R_P$ partial rounds, and the entire sequence
collapses, on $V$, to the linear map $\delta\mapsto M_{22}^{R_P}\delta$.
\end{lemma}
\begin{proof}
If $A_0=B_0$ then $A_0^{\alpha}=B_0^{\alpha}$, so the post-S-box
difference of two states differing only in $V$ is still $(0,0,\delta)$;
round-constant addition is affine and preserves this; by hypothesis
$M\cdot(0,0,\delta)\in V$, so the difference re-enters $V$. Induction
over $R_P$ rounds gives the claim.
\end{proof}

For $M=I$ ($B_d(I)=2<4$), a simulated diffusion trace over six partial
rounds at $\alpha=5$ shows the diffusion coefficient frozen at $1/3$
indefinitely, versus full diffusion after a single round under an MDS
benchmark matrix. MDS is thus necessary. We spend the rest of this paper
on why it is not sufficient.

\section{Related Work}
\label{sec:related}
Invariant-subspace analysis of stationary linear layers originates with
GRS21's sufficient condition for excluding infinitely long subspace
trails~\cite{weaklayers2020}, complemented by Ashur, Buschman, and
Mahzoun's broader algebraic cryptanalysis of the HADES
strategy~\cite{algebraiccrypt2023} and Keller and Rosemarin's
re-examination of HADES's middle layers~\cite{mindthemiddle2021}.
Li and Wang~\cite{cryptoeprint:2026/967} and, independently, Vitto
\cite{slipway2026} formalize that \emph{finite} subspace trails are
unavoidable for every linear layer, not themselves a flaw; Vitto's
contribution is showing these can be reached via full-round absorption
under adaptive matrix choice. Grassi, Koschatko, and Rechberger use
subspace trails to accelerate Gr\"obner-basis cryptanalysis across
operational modes~\cite{poseidonneptune2025}, and Li, Liu, and Wang
extend the trail concept to nonlinear subspaces~\cite{nonlinear2026}.
Merz and Rodr\'iguez Garc\'ia~\cite{cryptoeprint:2026/306} exhibit a
structurally distinct mechanism specific to Poseidon2(b)'s
Kronecker-structured external matrix. Jo~\cite{jo2026}, developed in
parallel with and independently of both, shows that the finite-trail and
tensor-structure mechanisms are not the only way a properly MDS,
certificate-clean matrix can be broken: an adaptively interpolated
matrix yields a complete compression collision. Jo's own discussion
names, without constructing or testing, the generation-order fix we
build and validate in \S\ref{sec:countermeasure}--\S\ref{sec:midpoint-reset}.
A comprehensive, independently authored 2026 survey of this entire
area~\cite{sok2026}, read in full for this paper, confirms (\S\ref{sec:sok})
that no existing work performs this synthesis or validation.

Beyond the permutation itself, Poseidon-based Fiat--Shamir protocols
have also been studied; \cite{cryptoeprint:2026/409} reports negative
results for the approaches it explores, and we do not extend any claim
in this paper to the protocol layer (\S\ref{sec:discussion}).
Bounty-driven resultant-based cryptanalysis of small-scale
instances~\cite{cryptoeprint:2026/150} and a generalized S-box-skipping
construction~\cite{sboxskip2026} are discussed in \S\ref{sec:discussion}
as attack classes our countermeasure does not address.

\section{Threat Model}
\label{sec:model}

\textbf{Setting.} The adversary knows the full public specification:
field, width, S-box exponent, round schedule, and (eventually) the
linear layer and round constants. We distinguish two generation
settings, matching the exact language the Poseidon Cryptanalysis
Initiative's own collision-challenge rules use~\cite{jo2026}:

\begin{itemize}[leftmargin=1.2em]
  \item \textbf{Non-adaptive (the setting this paper targets).} The
  complete parameter set -- both linear layers and every round constant
  -- is derived from one committed, public, non-adaptive transcript
  before any attack attempt begins. No party chooses any component with
  knowledge of another component that was not already fixed.
  \item \textbf{Adaptive (the setting all three attacks exploit).} Some
  component (typically the linear layer) may be selected with knowledge
  of another already-fixed component (typically the round constants),
  subject only to passing stated structural checks (MDS, subspace-trail
  clearance).
\end{itemize}

\textbf{Attack goal.} Given admissible inputs $\mathcal I$ and required
outputs $\mathcal O$, the constrained-input/constrained-output (CICO)
task is to find $x\in\mathcal I$ with $P(x)\in\mathcal O$, where $P$ is
the full permutation; a \emph{collision} additionally requires a second,
distinct admissible $x'$ with $H(x)=H(x')$ under the declared
feed-forward compression $H(x)=P(x)+x$. Jo's attack
(\S\ref{sec:midpoint-reset}) is a concrete instance of the second case.

\textbf{Scope.} We evaluate attacks and countermeasures at the
permutation/compression level only. We make no claim about any
higher-level protocol (e.g.\ a Fiat--Shamir transcript or a SNARK
verifier) built atop a Poseidon instance; \S\ref{sec:discussion}
discusses why such a claim would need separate, protocol-specific
argument.

\section{A Unified Structural Criterion}
\label{sec:criterion}

Both Vitto's and Merz and Rodr\'iguez Garc\'ia's mechanisms instantiate
the same underlying question, differing only in which matrix
($M_I$ vs.\ $M_{\mathrm{full}}$) and which round segment (partial vs.\
full) an unconstrained subspace survives through:

\begin{quote}
\emph{For a candidate $(M_I,M_{\mathrm{full}})$, does any proper subspace
of either matrix's eigenstructure survive into the nonlinear rounds
without being constrained by the generation procedure?}
\end{quote}

\begin{definition}[Finite-horizon kernel]
For stationary $M$ and active coordinate $0$, let
$K_r(M)=\bigcap_{j=0}^{r-1}\ker\big(e_0^{T}M^{j}\big)$, the subspace of
directions whose active-coordinate contribution is $X$-independent for
$r$ rounds, and $K_\infty(M)=\bigcap_{j\ge0}\ker(e_0^T M^j)$.
\end{definition}
\begin{proposition}
\label{prop:kinf}
$\dim K_r(M)\ge t-r$ for every $r$ (Vitto's unavoidability bound), and
$K_\infty(M)$ is the largest $M$-invariant subspace contained in
$\ker e_0^T$.
\end{proposition}
\begin{proof}
The first bound is rank--nullity on the $r\times t$ stack of rows
$e_0^TM^j$, $j<r$. For the second, Cayley--Hamilton expresses every
power $M^j$, $j\ge t$, as a combination of $I,\ldots,M^{t-1}$, so
$K_\infty=K_t$; if $v\in K_\infty$ then $e_0^TM^{j}(Mv)=e_0^TM^{j+1}v=0$
for all $j\ge0$, so $Mv\in K_\infty$ too, i.e.\ $K_\infty$ is
$M$-invariant; conversely any $M$-invariant $V\subseteq\ker e_0^T$
satisfies $e_0^TM^jV=0$ for every $j$, hence $V\subseteq K_\infty$.
\end{proof}

\begin{algorithm}[t]
\small
\caption{Unified structural certificate}
\label{alg:criterion}
\begin{algorithmic}[1]
\Require $M_I, M_{\mathrm{full}}\in\mathbb F_p^{t\times t}$, round-constant stream
\State \textbf{(a)} Factor $\chi(M_I)$ over $\mathbb F_p$; for each irreducible
factor $f_i$, compute $\ker f_i(M_I)$ -- the general GRS21 route
\State \textbf{(a-fast)} Check irreducibility of $\chi(M_I^{\,i})$ for
$i=1,\ldots,4(t-1)$: if irreducible for some $i$, $K_\infty(M_I)=0$
without (a)
\State \textbf{(b)} Compute $K_r(M_I)$ for $r=0,\ldots,t-1$; verify
$\dim K_r=t-r$ (``maximal rank growth,'' the tightest instance of
Proposition~\ref{prop:kinf} achievable)
\State \textbf{(c)} For every irreducible factor of $\chi(M_{\mathrm{full}})$,
compute its invariant subspace and flag it if its basis exhibits a
contiguous zero-run $\ge t/4$ -- generalizes~\cite{cryptoeprint:2026/306}'s
tensor-eigenvector finding to any $M_{\mathrm{full}}$, no Kronecker
structure assumed
\State \Return (a/a-fast) status, (b) profile, (c) flagged subspaces
\end{algorithmic}
\end{algorithm}

\paragraph{Implementation and validation.} Algorithm~\ref{alg:criterion}
is implemented field-generically (the prime is inferred from the data
handed to it, not hardcoded) and validated against exact controls before
being trusted on any candidate. For $M=I$, $t=3$: step (a-fast) correctly
fails and step (b) gives $\dim K_r=[3,2,2]$ for $r=0,1,2$ -- failing to
shrink at $r=2$, correctly flagging the known infinite trail of
Lemma~\ref{lem:invariant-subspace}. Against our own NUMS candidates
(\S\ref{sec:lightweight}) at $t=3$ and $t=16$: only the
\texttt{poseidon2\_sparse} family is irreducible at every tested power
(step a-fast succeeds); denser controls are secure only via the general
route (step a) -- a genuine empirical instance of ``GRS21 is sufficient,
not necessary,'' not an assumption. Step (b) gives exact equality
$\dim K_r(M_I)=t-r$ for every tested $r$ on our final candidate,
confirming its unavoidable finite trail is no longer than
Proposition~\ref{prop:kinf} guarantees. Applied to the real
$16\times16$ matrices $M_\varepsilon=P_4\otimes M_4$ and the swapped
$\bar M_\varepsilon=M_4\otimes P_4$, step (c) correctly separates them
(\texttt{suspicious=True} / \texttt{False} respectively) \emph{without
being told the Kronecker structure}.

\paragraph{What this criterion cannot see.} Every check above is a
static property of $M_I$ or $M_{\mathrm{full}}$ in isolation. Jo's
construction produces an $M$ with no exploitable static property --
it is MDS, step (a)/(a-fast) report it secure, and step (c) finds no
suspicious subspace -- because the vulnerability is a \emph{relation}
between $M$ and the already-fixed constants. No static certificate,
however refined, can close this gap; \S\ref{sec:countermeasure} removes
it by controlling generation order instead.

\section{The Generation-Order Countermeasure}
\label{sec:countermeasure}

We bind the round constants to a canonical encoding of \emph{both}
linear layers before either is used:
\[
  T=\texttt{seed}\,\|\,\mathrm{enc}(p,t,\alpha,R_F,R_P)\,\|\,
  \mathrm{enc}(M_I)\,\|\,\mathrm{enc}(M_{\mathrm{full}}),
\]
\[
  c^{(r)}_i = \mathrm{HashToField}\big(T\,\|\,\mathrm{enc}(r,i)\big).
\]
Both matrices are generated first, non-adaptively, via a
nothing-up-my-sleeve (NUMS) procedure: a domain-separated SHAKE-256
hash-to-field, rejection-sampled for invertibility (and, for one
candidate family, for Algorithm~\ref{alg:criterion} step (a-fast)
irreducibility). Round constants are then deterministic functions of a
transcript that already contains both matrices. No constant can be fixed
independently of either matrix, and no matrix can be chosen as a
function of constants that do not yet exist. This removes, by
construction, the precondition Vitto's and Jo's attacks both require,
and the same precondition the kleptographic ``nothing-up-my-matrix''
backdoor class needs to hide a weakness behind an otherwise-passing
matrix~\cite{backdoor2026}.

\textbf{A real gap, found and closed in adversarial self-review.} An
earlier version of this binding covered only $M_I$. Vitto's own attack
constructs the \emph{full-round} matrix adversarially against fixed
constants (absorption into $M_I$'s trail happens through the full-round
\emph{prefix}, which $M_{\mathrm{full}}$ governs) -- binding only $M_I$
would leave exactly this door open in any standards process that lets a
submitter choose their own external matrix, which is precisely Vitto's
own target setting. We fixed this before any result below was produced.

\textbf{The swapped Kronecker order (closes~\cite{cryptoeprint:2026/306}).}
For a tensor-structured $M_{\mathrm{full}}=M_4\otimes P_{t/4}$ (the
factor order swapped from the vulnerable $P_{t/4}\otimes M_4$), identical
multiplicative cost destroys the exploitable $(q,-q,0,\ldots,0)$ shape:
confirmed both by Algorithm~\ref{alg:criterion} step (c) (zero suspicious
contiguous zero-runs) and by a dynamic check pushing a generic basis
vector through one real full round at both BN254 and KoalaBear scale.
This is not a paper-only concern: reading Plonky3's production
\texttt{mds\_light\_permutation} code directly shows it uses the
\emph{unswapped}, vulnerable order in an actively used STARK proving
stack (\S\ref{sec:audit}).

\section{Replicating and Defeating Jo's Full-Round Collision}
\label{sec:midpoint-reset}

Jo's discussion proposes, without building or testing it, exactly
\S\ref{sec:countermeasure}'s binding: \emph{``derive the complete
parameter set non-adaptively\dots by binding the round constants to a
canonical encoding of the matrix.''} This section reports, to our
knowledge, the first construction and empirical test of that claim.

\subsection{The attack mechanism}
Track two executions' midpoint $m_r=(x^+_r+x^-_r)/2$ and half-difference
$d_r=(x^+_r-x^-_r)/2$. One prescribed matrix image per two-round block,
$Mu_r=-c^{(r+1)}$, cancels the midpoint against the next round constant;
since the S-box is odd, the following S-box's output midpoint is exactly
zero and stays zero for the rest of the permutation (``midpoint
reset''). Two reusable images, $Me=v$, $Mw=e$ on a coordinate $e$
passive in every partial round, keep the half-difference confined to
$\mathbb F_p\cdot e$ through every full \emph{and} partial round.
Collecting these images as columns of $U$ (sources, computable from the
public constants alone, not from $M$) and $V$ (targets), $M=VU^{-1}$ is
the unique matrix realizing the whole trajectory. A correctly closed
scalar trajectory (final difference the negative of the initial
difference) yields a genuine feed-forward collision in an $M$ that is
MDS and clears every certificate of \S\ref{sec:criterion}.

\subsection{Replication methodology}
We implemented the construction generically (any $t$, field, round
constants), with one deliberate scope reduction relative to Jo's own
$(t,\alpha,R_F,R_P)=(16,3,8,20)$ KoalaBear instance: $R_F=4$ (one
full-round block per side, instead of two), which collapses the harder
two-block cubic entry/exit solve to a closed-form single step. Choosing
the entry block's free scalar $s_0=0$ gives $b_0=\sqrt[3]{a}$ directly
(well-defined since $\gcd(3,p-1)=1$, the condition $\alpha=3$ already
requires), and the block's output lands at exactly $a$ by construction.
The one remaining free scalar -- the closing amplitude $a$ -- is found
by brute-force search over a small toy prime field, the same
replicate-the-mechanism-not-the-literal-instance discipline used
throughout this line of work for Vitto's own $3^{10}$ vs.\ $3^{28}$
numbers.

\textbf{Result.} Against unprotected (unbound) constants, at
$t=6,\alpha=3,R_F=4,R_P=4$ over a toy field $p=10037$, closing amplitude
$a^{*}=177$ yields a genuine, independently-verified feed-forward
collision,
\[
  H(x^+)=H(x^-)=(0,\,3518,\,0,\,0,\,0,\,0),\qquad x^+\ne x^-,
\]
confirmed by executing the actual permutation forward with the solved
$M$, not merely the construction's own scalar bookkeeping. One real
implementation bug was caught in the process: an early version computed
$M=U^{-1}V$ instead of $M=VU^{-1}$ (matrix multiplication is
non-commutative, silently solving a different system) -- caught by a
direct $M\,U\stackrel{?}{=}V$ equality check, which we now retain
permanently as a verification step on every solved candidate.

\subsection{Immunity test}
We generated a matrix $M$ the honest way
(\S\ref{sec:countermeasure}), derived constants
$c=\mathrm{Hash}(\texttt{seed}\,\|\,\mathrm{params}\,\|\,M\,\|\,M)$ --
constants fixed \emph{after and from} $M$, the reverse of the order the
attack needs -- and scanned \emph{every} nonzero $a$ in the toy field
for whether this already-fixed $M$ satisfies the attack's required image
constraints $M\,U(c,a)=V(c,a)$ for any of them.

\begin{table}[t]
\centering
\caption{Binding-immunity test: consistent closing amplitudes found out
of $p-1$ candidates scanned, across independent trials.}
\label{tab:immunity}
\small
\begin{tabular}{@{}llcc@{}}
\toprule
Trial set & $M_I$ family & Scanned & Found \\
\midrule
5 seeds, $p=10007$ & \texttt{poseidon2\_sparse} & 10006 ea. & 0 \\
$p=10007$ & \texttt{low\_weight\_circulant} & 10006 & 0 \\
$p=10007$ & \texttt{irreducible\_filtered} & 10006 & 0 \\
$p=10067$ & \texttt{low\_weight\_circulant} & 10066 & 0 \\
$p=10067$ & \texttt{irreducible\_filtered} & 10066 & 0 \\
\bottomrule
\end{tabular}
\end{table}

Zero consistent amplitudes were found in all nine independent trials
(Table~\ref{tab:immunity}). We state the result at the confidence it
supports: this is empirical evidence, not a formal theorem, that
honestly NUMS-generated, bound $(M,c)$ pairs do not satisfy the attack's
highly constrained algebraic relation -- consistent with the heuristic
that an $M$ uncorrelated with the attack's image constraints by
construction satisfies them only with probability on the order of the
field size, and with Jo's own qualitative argument that binding removes
the ``choose $M$ after $c$'' freedom the construction depends on. A
formal fixed-point/preimage-hardness argument, and replication at Jo's
literal $R_F=8$/KoalaBear scale, are future work
(\S\ref{sec:discussion}).

\section{Lightweight Instantiation}
\label{sec:lightweight}

Candidate $M_I$ matrices are generated via three NUMS families:
\texttt{poseidon2\_sparse} (near-identity off-diagonal, hash-derived
diagonal -- the real Poseidon2 internal-matrix shape),
\texttt{low\_weight\_circulant}, and \texttt{irreducible\_filtered}
(dense, rejection-sampled for Algorithm~\ref{alg:criterion} step
(a-fast)). We separate two cost measures the literature sometimes
conflates:
\[
  \mathrm{cost}_{\mathrm{R1CS}}(t,R_F,R_P,\alpha)=R_F\cdot t\cdot
  c_\alpha+R_P\cdot c_\alpha,\quad
  \mathrm{mul}(M)=|\{M_{ij}\notin\{0,1\}\}|,
\]
with $c_\alpha$ the square-and-multiply constraint count for
$x^\alpha$ ($c_5=3,c_3=2$); in a standard R1CS encoding, multiplying a
wire by a public constant is a free linear combination.

\begin{table}[t]
\centering
\caption{Candidate comparison at $t=3,\alpha=5,R_F=8,R_P=57$ (BN254).}
\label{tab:candidates}
\small
\begin{tabular}{@{}lccc@{}}
\toprule
Family & R1CS n.l. & R1CS lin. & $\mathrm{mul}(M_I)$ \\
\midrule
\texttt{poseidon2\_sparse} & 243 & 276 & 195 \\
\texttt{low\_weight\_circulant} & 243 & 276 & 390 \\
\texttt{irreducible\_filtered} & 243 & 276 & 585 \\
\bottomrule
\end{tabular}
\end{table}

R1CS cost is \textbf{identical} across all three candidates (confirmed by
an independent \texttt{circom} compile, not just the formula): a
circuit row is allocated per matrix \emph{row}, not per nonzero entry.
Sparsity is a real win for native/Plonkish multiplicative complexity and
has zero measurable R1CS effect. \texttt{poseidon2\_sparse} wins on the
metric that differs and is used for the remainder of this paper.

\textbf{Cross-validation against production.} Compiling our
NUMS-generated circuit and \texttt{iden3/circomlib}'s production
\texttt{Poseidon(2)} template at identical parameters gives
byte-identical non-linear constraint counts (243 for both) and
near-identical linear counts (276 vs.\ 274, cosmetic hash-wrapper
exposure). A full Groth16 pipeline (Powers-of-Tau, setup, prove, verify)
produced valid, independently-verified proofs for both.

\textbf{Field size.} At $t=3,\alpha=3$ over KoalaBear, a first attempt
reused $R_F=8$ from the public $t=16$ instance; GKR+21's own statistical
bound~\cite{poseidon2021} (Eq.~2, scaling with $t+1$) requires $R_F=12$
at $t=3$. Corrected: $112$ non-linear $+152$ linear $=264$ total
constraints, vs.\ $512$--$519$ at BN254 -- a real
$\approx1.9\times$ reduction, confirmed by
\texttt{circom --prime goldilocks}; a deployment decision, not an
algorithmic improvement.

\textbf{Tensor-structuring $M_{\mathrm{full}}$.} At $t=16$ with the safe
factor order: identical R1CS cost ($892$ total, byte-identical to an
all-dense control) but a $73\%$ reduction in real multiplicative
complexity ($640$ vs.\ $2368$) once costed via the real fast block
algorithm (apply $M_4$ to each of $t/4$ blocks, then a pure-addition
outer mix) rather than a naive entry count, which misses the difference
by $4.6\times$.

\subsection{Negative results}
\label{sec:levers}
We report, rather than omit, five tested levers that did not help.
\textbf{(1) Lowering $\alpha$.} At $t=3$, $\alpha{:}5{\to}3$ requires
$R_P{:}56{\to}83$; the extra linear-constraint overhead exceeds the
S-box saving ($594$ vs.\ $512$ total). \emph{Heavier, not lighter.}
\textbf{(2) Rotating the active S-box coordinate.} Offers no benefit for
matrices already at Proposition~\ref{prop:kinf}'s dimension-decay
ceiling, and \emph{introduces} a new infinite trail in a matrix
(\texttt{circulant}) secure under the fixed-position schedule.
\emph{Double-edged.} \textbf{(3)--(4) Two GMiMCHash2 tricks.} The
$\alpha=2^k$ S-box does not port: GMiMC is Feistel (invertible without a
bijective round function); our SPN structure requires
$\gcd(\alpha,p-1)=1$, which no power of two satisfies over an
odd-characteristic field. The $M_{IO}$ input/output matrix targets proof
tractability, not cost, and moves no cost number. \textbf{(5) A naive
lookup-hybrid model.} Modeling a ``lookup round'' as near-maximal
algebraic degree produces a vacuous result -- replacing 56 of 57 partial
rounds with one lookup round ``clears'' the model because it cannot
represent an attacker simply attacking the one remaining algebraic round
separately. \emph{Invalidated, not negative evidence either way.}

\section{Exhaustive Validation}
\label{sec:battery}

The final candidate ($t=16,\alpha=3,R_F=8,R_P=20$, KoalaBear,
$M_I=\texttt{poseidon2\_sparse}$, $M_{\mathrm{full}}=M_4\otimes P_4$
safe order, constants bound to both) was run through every available
check.

\begin{table}[t]
\centering
\caption{Exhaustive battery against the final candidate.}
\label{tab:battery}
\small
\begin{tabularx}{\columnwidth}{@{}p{0.04\columnwidth}Xp{0.4\columnwidth}@{}}
\toprule
\# & Check & Result \\
\midrule
1 & Statistical round sufficiency (Eq.~2,~\cite{poseidon2021}) & $\min
R_F=8$ (bound 6, margin 2); candidate $R_F=8$ sufficient, at the
boundary \\
2 & MDS / branch number of $M_I$ & $B_d(M_I)=17=t+1$, no slack \\
3 & Alg.~\ref{alg:criterion}(a/a-fast) & \texttt{SECURE} \\
4 & Alg.~\ref{alg:criterion}(c), $M_{\mathrm{full}}$ & flags one
subspace; investigated, confirmed a \textbf{false positive} (resolved
by item 7) \\
5 & Alg.~\ref{alg:criterion}(b) & $\dim K_r=t-r$, $r=0,\ldots,15$ --
maximal rank growth \\
6 & Symbolic freeze, $d\in K_4(M_I)$ & $[0,0,0,0,1,3]$ -- frozen 4
rounds, jumps at round 4 exactly \\
7 & No-accidental-landing, generic direction, $M_{\mathrm{full}}$ &
$[1,3,9]=\alpha^{\text{round}}$ exactly -- no accidental trail \\
8 & Real multiplicative cost & $640$ \\
\bottomrule
\end{tabularx}
\end{table}

\textbf{Item 4, resolved -- a concrete certificate-scope illustration.}
The static scan flags the \emph{safe}-order $\bar M_\varepsilon$ at
KoalaBear, clean at BN254 -- field-dependent factorization, not a
contradiction. A generic (non-$\pm1$) basis vector of the flagged
subspace reaches full weight after one real full round: no
negation-compatible structure survives, confirming a genuine-but-inert
subspace. Lesson: a static flag must always be paired with a dynamic
one-round survival check before being reported as a finding.

\subsection{Real-world audit}
\label{sec:audit}
We ran Algorithm~\ref{alg:criterion} against three deployed parameter
sets we did not generate: \textbf{\texttt{circomlib}}'s real $t=3$
matrix is clean on every check. \textbf{Plonky3}'s round numbers match
our statistically-derived formula and~\cite{sok2026}'s recomputed table
exactly; its external matrix uses the vulnerable, unswapped Kronecker
order, confirmed by code reading, reproducing the matrix, and our
Kronecker-agnostic certificate flagging it unprompted -- practical
exploitability depends on a rate/capacity parameter set per call site,
which we did not trace. \textbf{\texttt{arkworks-rs}} ships no fixed
matrix; its Cauchy-matrix generator's own source comments claim three
well-formedness checks with zero corresponding code -- a documented
implementation gap, not a demonstrated break of a specific deployed
instance.

\section{Positioning Against the State of the Art}
\label{sec:sok}

A full read of~\cite{sok2026} (134 pages, read in full after earlier
access attempts were blocked) gives real, already-computed benchmark
numbers we cite rather than re-derive.

\begin{table}[t]
\centering
\caption{R1CS constraints, single permutation, Gnark/Groth16, BLS12-381,
$t=4$ (\cite{sok2026}, Table 23).}
\label{tab:sok}
\small
\begin{tabular}{@{}lc@{}}
\toprule
Construction & R1CS constraints \\
\midrule
Griffin & 224 \\
Anemoi~\cite{anemoi2022} & 244 \\
Arion & 254 \\
Poseidon~\cite{poseidon2021} & 268 \\
Poseidon2~\cite{poseidon22023} & 268 \\
Rescue-Prime & 368 \\
Skyscraper & 1425 \\
Polocolo~\cite{polocolo2025} & 1553 \\
\bottomrule
\end{tabular}
\end{table}

Two findings bound this paper's novelty claims precisely.
\cite{sok2026}'s cryptanalysis treatment mentions both Vitto's and Jo's
attacks in one combined sentence, with no mechanism detail, no
countermeasure, and no round-number consequence -- confirming
\S\ref{sec:criterion}--\S\ref{sec:midpoint-reset} are not duplicated
anywhere in the literature. \cite{sok2026} also confirms ``no instances
require updating under the latest attacks'' for Poseidon/Poseidon2's
round-number formula, independently validating item~1 of
Table~\ref{tab:battery}.

\section{Discussion and Limitations}
\label{sec:discussion}

\textbf{Certificate scope.} A zero-kernel or clean-scan result excludes
only the specified static model (\S\ref{sec:criterion}); it is not
evidence against generation-order attacks. \S\ref{sec:battery}'s item-4
false positive and \S\ref{sec:midpoint-reset}'s central finding (a
matrix clearing every certificate in this paper while still admitting a
complete collision) are the two concrete illustrations this paper's own
results produced of exactly this limit.

\textbf{Protocol-layer limits.} Every claim is restricted to the
permutation/compression task evaluated. Small-instance successes
\cite{cryptoeprint:2026/150} do not extrapolate to production
parameters by themselves, and a permutation weakness does not establish
protocol forgery; \cite{cryptoeprint:2026/409}'s own negative results
for the Fiat--Shamir layer it explores support restraint here, not an
automatic conclusion either way.

\textbf{Open problems.} Two attack classes are outside this paper's
countermeasure, confirmed genuinely open by a full read of the most
comprehensive survey available (\S\ref{sec:sok}), not merely by our own
prior gap: \textbf{S-box skipping}~\cite{sboxskip2026}, claimed
independent of matrix, constants, exponent, and field, is cited only in
passing by~\cite{sok2026} with no mechanism analysis; and
\textbf{bivariate resultant attacks against Poseidon's own linear
layer}~\cite{resultant2026} -- the strongest published results in this
family target Type-II primitives (Rescue-Prime, Anemoi, Griffin, Arion),
not Poseidon. A \textbf{formal fixed-point/preimage-hardness argument}
for \S\ref{sec:countermeasure}'s binding, beyond the nine-trial
empirical test, and \textbf{replication at Jo's literal
$R_F=8$/KoalaBear scale} (needing a general cubic-root solver rather
than our closed-form $R_F=4$ shortcut), are also left open.

\section{Ethics Considerations}
Every attack this paper analyzes was previously published by its
original authors; we discovered no new vulnerability in any deployed
system. Our real-world audit (\S\ref{sec:audit}) identifies a
structural precondition already disclosed in the published
literature~\cite{cryptoeprint:2026/306} (Plonky3's Kronecker factor
order) and a documentation/implementation gap
(\texttt{arkworks-rs}'s unchecked generator); we did not construct a
working exploit against any production deployment, and we state
explicitly that Plonky3's practical exploitability depends on a
rate/capacity parameter we did not determine. We release this analysis
to help deployers and standards processes adopt the generation-order
binding of \S\ref{sec:countermeasure} before, not after, an adaptive-matrix
attack is demonstrated against a system they operate.

\section{Open Science}
The complete implementation (field-generic linear algebra, the
unified-criterion scanner, the NUMS candidate generator, the Midpoint
Reset replication and immunity test, and the \texttt{circom}/Groth16
benchmarking harness), all round-constant and matrix generation seeds,
and the full raw experimental log underlying every number reported in
this paper are maintained in a version-controlled repository, including
failed attempts and the bugs found and fixed in the course of this work.
We intend to make this repository publicly available upon publication.

\section{Conclusion}
MDS is necessary for Poseidon's partial-round linear layer, but three
independent 2026 results show it is not sufficient: properly MDS,
certificate-clean matrices can still be broken via finite-trail
absorption, tensor-structured round skipping, or an adaptively chosen
matrix programming an exact collision trajectory. We identified the one
assumption all three violate -- that the linear layer and round
constants are generated independently -- gave a unified structural
criterion for everything that assumption's violation does \emph{not}
explain, and a single generation-order countermeasure for everything it
does. We validated that countermeasure empirically against all three
attacks, including, to our knowledge, the first replication and defense
validation of Jo's full-round collision reported anywhere, and
instantiated the result as a lightweight, production-cross-validated
KoalaBear candidate. S-box skipping and bivariate resultant attacks
against Poseidon's own linear layer, and a formal hardness argument for
our countermeasure, remain open problems.

\printbibliography

\end{document}
